% BHU PYQ -- Transcribed & Corrected
% Paper: CHEMJ44: Techniques in Chemistry
% Exam: B.Sc. (Hons) Semester IV Examination 2025-26 (NEP)
% Ref Code: Conf/Even/N/2025-26/09
% Category: Major
% Department: Chemistry

\documentclass[11pt,a4paper]{article}
\usepackage[a4paper,top=2.5cm,bottom=2.5cm,left=2.8cm,right=2.8cm,headheight=14pt]{geometry}
\usepackage{amsmath,amssymb}
\usepackage[T1]{fontenc}
\usepackage[utf8]{inputenc}
\usepackage{lmodern,microtype,fancyhdr,enumitem,hyperref,lastpage,parskip}

\hypersetup{
    colorlinks=true,
    urlcolor=black,
    linkcolor=black,
    pdftitle={CHEMJ44: Techniques in Chemistry -- BHU Sem IV 2025-26 (NEP)}
}

\pagestyle{fancy}
\fancyhf{}
\renewcommand{\headrulewidth}{0.4pt}
\renewcommand{\footrulewidth}{0.4pt}
\fancyhead[L]{\small \textit{Banaras Hindu University}}
\fancyhead[R]{\small \textit{CHEMJ44: Techniques in Chemistry (2025-26)}}
\fancyfoot[C]{\small Page \thepage\ of \pageref{LastPage}}

\newcommand{\pts}[1]{\hfill \textit{[#1]}}

\begin{document}

\begin{center}
  {\large\bfseries BANARAS HINDU UNIVERSITY}\\[2pt]
  {\normalsize B.Sc. (Hons) Semester IV Examination 2025-26 (NEP)}\\[6pt]
  \rule{0.8\linewidth}{0.5pt}\\[6pt]
  {\large\bfseries CHEMISTRY}\\[4pt]
  {\large\bfseries Paper Code: CHEMJ44 : Techniques in Chemistry}\\[2pt]
  {\small Ref. No.: Conf/Even/N/2025-26/09}\\[4pt]
  {\normalsize Time: 2$\frac{1}{2}$ Hours\hfill Maximum Marks: 70}
\end{center}

\rule{\linewidth}{0.4pt}

\smallskip

\noindent \textit{Instructions: Answer five questions including Question 1 which is compulsory. The figures in the right-hand margin indicate marks.}

\bigskip

\noindent \textbf{Question 1.}
\begin{enumerate}[label=(\alph*)]
  \item What is a standard solution? Write the difference between primary and secondary standard solutions with suitable examples. \pts{2}
  \item What is an indicator? Describe its different types with suitable examples. \pts{2}
  \item Sketch the various types of bending vibration in IR of $\text{H}_2\text{O}$. \pts{2}
  \item Why pH ranges from 0 to 14 only? \pts{2}
  \item Calculate the pH value of $10^{-9}\text{ mole HCl}$. \pts{2}
  \item What is crystallisation? State its importance in purification. \pts{2}
  \item What is the difference between qualitative and quantitative analysis? Explain with examples. \pts{2}
\end{enumerate}

\bigskip

\noindent \textbf{Question 2.}
\begin{enumerate}[label=(\alph*)]
  \item Describe Instrumentation of UV-Visible absorption spectroscopy and its applications. \pts{4}
  \item Explain the concepts of molarity, normality and molality. Derive their expressions and discuss their interrelationships with suitable numerical examples. \pts{5}
  \item Explain the principle of complexometric titration. Describe the role of EDTA, metal ion indicators and applications of this method. \pts{5}
\end{enumerate}

\bigskip

\noindent \textbf{Question 3.}
\begin{enumerate}[label=(\alph*)]
  \item Discuss the process of crystallization and precipitation. Explain the factors affecting crystal formation, purity and yield, along with their applications in chemical purification. \pts{6}
  \item Discuss the concept of conjugate acid-base pairs. \pts{2}
  \item Explain the general concepts of acids and bases according to Arrhenius, Brønsted-Lowry and Lewis theories, with suitable examples. \pts{6}
\end{enumerate}

\bigskip

\noindent \textbf{Question 4.}
\begin{enumerate}[label=(\alph*)]
  \item Define the following terms used in UV-Visible absorption spectroscopy, with suitable examples: \pts{6}
  \begin{enumerate}[label=(\roman*)]
    \item Bathochromic shift
    \item Hypsochromic shift
    \item Hypochromic effect
    \item Hyperchromic effect
  \end{enumerate}
  \item Discuss the principles and procedures of acid-base titrations. Explain the titration curves for a strong acid vs. a strong base and a weak base vs. a strong acid. \pts{4}
  \item A $25\ \mu\text{L}$ serum sample was analyzed for glucose content and found to contain $26.7\ \mu\text{g}$. Calculate the concentration of glucose in $\mu\text{g/ml}$. \pts{4}
\end{enumerate}

\bigskip

\noindent \textbf{Question 5.}
\begin{enumerate}[label=(\alph*)]
  \item Describe the principle of chromatography. Compare paper chromatography, thin-layer chromatography (TLC) and column chromatography in terms of method, advantages and applications. \pts{8}
  \item What is the Lambert-Beer law? Explain each term in the expression. \pts{2}
  \item Write short note on chemical shift in NMR. \pts{4}
\end{enumerate}

\bigskip

\noindent \textbf{Question 6.}
\begin{enumerate}[label=(\alph*)]
  \item Describe infrared (IR) spectroscopy. Explain different IR sources. How we prepare sample for analysis? \pts{5}
  \item Describe redox titrations, including their principles, types and applications. \pts{5}
  \item Write notes on pH and pH calculations. \pts{4}
\end{enumerate}

\bigskip

\noindent \textbf{Question 7.}
\begin{enumerate}[label=(\alph*)]
  \item Explain the principles of distillation. Describe in detail steam distillation and vacuum distillation, including their working mechanisms, advantages and applications. \pts{8}
  \item Differentiate between the fingerprint region and the group frequency region in IR spectroscopy. \pts{3}
  \item Explain precipitation titrations in details. \pts{3}
\end{enumerate}

\bigskip
\begin{center}
  \textbf{***}
\end{center}

\end{document}
